
\chapter{Geometry in $\mathbb{R}^2$ and $\mathbb{R}^3$}

\section{Describing Lines in $\mathbb{R}^2$ using Vectors}

\subsection{Homogeneous and Nonhomogeneous Equations}

A line in $\mathbb{R}^2$ is given algebraically by a single linear equation.
The simplest lines to study are also those that contain the origin $(0,0)$.

$$
\text{Example: Let $\ell$ be the line}
\quad 
3x-4y=0.
$$

\begin{figure}[h]
\centering
\begin{tikzpicture}[scale=0.72]
  % Axes
  \draw[->, gray] (-5.5,0) -- (5.8,0) node[right] {$x$};
  \draw[->, gray] (0,-4.4) -- (0,4.8) node[above] {$y$};

  % Line 3x - 4y = 0
  \draw[black, thick, domain=-5.2:5.2] plot (\x,{0.75*\x});

  % Points
  \filldraw[black] (0,0) circle (2pt);
  \node[below right] at (0,0) {$(0,0)$};

  \filldraw[black] (4,3) circle (2pt);
  \node[right=8pt] at (4,3) {$(4,3)$};

  \filldraw[black] (-4,-3) circle (2pt);
  \node[below left] at (-4,-3) {$(-4,-3)$};

  % Label for the line
  \node[violet, above left] at (5.05,4.05) {$\ell: 3x-4y=0$};
\end{tikzpicture}
\caption{The line $3x-4y=0$, which passes through the origin.}
\end{figure}

\paragraph{Remark.}
The $0$ on the right of the equation is what forces the line to go through
$(0,0)$. Such an equation is called homogeneous.

If one changes the $0$ to a nonzero number, then the line will not go through
the origin $(0,0)$.

$$
\text{Example:}
\qquad
3x-4y=7
\qquad
\longrightarrow
\qquad
\text{nonhomogeneous linear equation.}
$$

\begin{figure}[h]
\centering
\begin{tikzpicture}[scale=0.72]
  % Axes
  \draw[->, gray] (-5.5,0) -- (6.5,0) node[right] {$x$};
  \draw[->, gray] (0,-4.5) -- (0,4.8) node[above] {$y$};

  % Homogeneous line 3x - 4y = 0
  \draw[black, thick, domain=-5.2:5.2] plot (\x,{0.75*\x});

  % Affine/nonhomogeneous line 3x - 4y = 7
  % y = (3x - 7)/4
  \draw[magenta!80!black, thick, domain=-1.2:6.2] plot (\x,{(3*\x - 7)/4});

  % Points on homogeneous line
  \filldraw[black] (0,0) circle (2pt);
  \node[below right] at (0,0) {$(0,0)$};

  \filldraw[black] (4,3) circle (2pt);
  \node[right=8pt] at (4,3) {$(4,3)$};

  \filldraw[black] (-4,-3) circle (2pt);
  \node[below left] at (-4,-3) {$(-4,-3)$};

  % Points on affine line
  \filldraw[magenta!80!black] (1,-1) circle (2pt);
  \node[below right=3pt, magenta!80!black] at (1,-1) {$(1,-1)$};

  \filldraw[magenta!80!black] (5,2) circle (2pt);
  \node[right=8pt, magenta!80!black] at (5,2) {$(5,2)$};

  % Arrow along affine line from (1,-1) to (5,2)
  
  % Labels
  \node[violet, above left] at (5.05,4.05) {$\ell: 3x-4y=0$};
  
  \node[magenta!80!black, above right] at (6.5,2.75) {$\ell': 3x-4y=7$};
  
\end{tikzpicture}
\caption{Changing the right-hand side from $0$ to $7$ gives a parallel line.}
\end{figure}

We get in this way a line that is parallel to $3x-4y=0$.

By varying $c$, the equation

$$
3x-4y=c
$$

will give all possible lines in the plane that are parallel to the line $\ell$.

For simplicity, we'll first study extensively the homogeneous case, i.e., lines
through the origin.

\subsection{Describing Lines Through the Origin Using Vectors}

Lines through the origin can be efficiently described using vectors.

Let

$$
\vec v=
\begin{bmatrix}
4\\
3
\end{bmatrix}.
$$

The line $\ell$ can be described as all multiples of this vector.


Indeed every point on the line is of the form $(3t,4t)$ for some scalar $t \in \mathbb{R}$. For example, we've plotted the vectors $\vec{v}$, $\frac{1}{2}\vec{v}$ and 
$-\frac{1}{2}\vec{v}$
in the figure below.  

\bigidea{We say $\ell$ is the span of the vector
$$
\vec{v} = \begin{bmatrix}
4\\
3
\end{bmatrix}.
$$
}



\begin{figure}[h]
\centering
\begin{tikzpicture}[scale=0.72]
  % Axes
  \draw[->, gray] (-5.5,0) -- (5.8,0) node[right] {$x$};
  \draw[->, gray] (0,-4.4) -- (0,4.8) node[above] {$y$};

  % Line
  \draw[black, thick, domain=-5.2:5.2] plot (\x,{0.75*\x});

  % Points
  \filldraw[black] (0,0) circle (2pt);
  \node[below right] at (0,0) {$(0,0)$};

  \filldraw[black] (4,3) circle (2pt);
  \node[right=8pt] at (4,3) {$(4,3)$};

  \filldraw[black] (-4,-3) circle (2pt);
  \node[below left] at (-4,-3) {$(-4,-3)$};

  % Vector arrow from origin to v = (4,3)
  \draw[-{Latex[length=3mm]}, thick] (0,0) -- (4,3);
  \node[above left=2pt] at (3.75,2.85) {$\vec v$};

  % Vector arrow from origin to (1/3)v = (4/3,1)
  \draw[-{Latex[length=2mm]}, thick] (0,0) -- (1.333,1);
  \filldraw[black] (1.333,1) circle (1.6pt);
  \node[above left=3pt] at (1.333,1) {$\frac{1}{3}\vec v$};

  % Vector arrow from origin to (1/2)v = (2,3/2)
  \draw[-{Latex[length=2.2mm]}, thick] (0,0) -- (2,1.5);
  \filldraw[black] (2,1.5) circle (1.6pt);
  \node[below right=3pt] at (2,1.5) {$\frac{1}{2}\vec v$};

  % Vector arrow from origin to -(1/2)v = (-2,-3/2)
  \draw[-{Latex[length=2.2mm]}, thick] (0,0) -- (-2,-1.5);
  \filldraw[black] (-2,-1.5) circle (1.6pt);
  \node[above left=3pt] at (-2,-1.5) {$-\frac{1}{2}\vec v$};

  % Labels
  \node[violet, above left] at (5.05,4.05) {$3x-4y=0$};

  \node[right] at (5.0,3.75) {$\ell$};
\end{tikzpicture}
\caption{The line $\ell$ through the origin can be described using the vector $\vec v=\begin{bmatrix}4\\3\end{bmatrix}$.}
\end{figure}








Note that we could have picked any nonzero vector along this line, e.g.

$$
\frac{1}{2}\vec{v} = \begin{bmatrix}
2\\
3/2
\end{bmatrix} \quad \text{or} \quad -\frac{1}{2}\vec{v} = \begin{bmatrix}
-2\\
-3/2
\end{bmatrix} \quad \text{or} \quad \frac{1}{3}\vec{v} = \begin{bmatrix}
4/3\\
1
\end{bmatrix}
$$
So in this case

$$
\ell=\operatorname{span}(\vec v)=\operatorname{span}(\frac{1}{2}\vec v)
=\operatorname{span}(-\frac{1}{2}\vec v)=\operatorname{span}(\frac{1}{3}\vec v) \quad \text{etc.}$$


\paragraph{Exercise.}
Let $\ell'$ be the span of

$$
\begin{bmatrix}
2\\
1
\end{bmatrix}.
$$

Draw $\ell'$. Give an equation describing $\ell'$.

\begin{figure}[h]
\centering
\begin{tikzpicture}[scale=0.72]
  % Axes
  \draw[->, gray] (-5.5,0) -- (6.5,0) node[right] {$x$};
  \draw[->, gray] (0,-3.2) -- (0,4.0) node[above] {$y$};

  % Line through origin with direction (2,1): y = x/2
  \draw[black, thick, domain=-5.2:6.2] plot (\x,{0.5*\x});

  % Origin
  \filldraw[black] (0,0) circle (2pt);

  % Vectors
  \draw[-{Latex[length=3mm]}, thick] (0,0) -- (2,1);
  \node[above left=4pt] at (2,1) {$\begin{bmatrix}2\\1\end{bmatrix}$};

  \draw[-{Latex[length=3mm]}, thick] (0,0) -- (6,3);
  \node[above=6pt] at (5.7,3) {$\begin{bmatrix}6\\3\end{bmatrix}$};

  % Label
  \node[right=6pt] at (5.4,2.7) {$\ell'$};
\end{tikzpicture}
\caption{The line $\ell'$ spanned by $\begin{bmatrix}2\\1\end{bmatrix}$.}
\end{figure}

Notice that the $y$-coordinate of any point on this line is always one half of
the $x$-coordinate. So the equation of the line is

$$
y-\frac12 x=0
$$

or

$$
x-2y=0.
$$

\paragraph{Remark.}
We could multiply the equation by any nonzero number, and that would give an
equally good equation.

For example,

$$
2x-4y=0.
$$

So the equation describing the line is not unique.

\subsection{The Dual Perspective and Orthogonality}

Let's reconsider the line

$$
\ell: 3x-4y=0.
$$

Let $\ell^\perp$ be the unique line through the origin that is perpendicular to $\ell$.


\begin{figure}[h]
\centering
\begin{tikzpicture}[scale=0.68]
  % Axes
  \draw[->, gray] (-5.8,0) -- (6.0,0) node[right] {$x$};
  \draw[->, gray] (0,-5.2) -- (0,5.3) node[above] {$y$};

  % Line ell: y = 3x/4
  \draw[black, thick, domain=-5:5.2] plot (\x,{0.75*\x});
  \node[right=5pt] at (4.5,4.45) {$\ell$};

  % Perpendicular line ell^\perp: direction (3,-4), slope -4/3
  \draw[black, thick, domain=-3.6:3.6] plot (\x,{-1.333333333*\x});
  \node[left=6pt] at (-2.6,5.47) {$\ell^\perp$};

  % Right angle marker
  \draw[thick]
    (0.36,0.27) --
    (0.09,0.63) --
    (-0.27,0.36);

  % Vector v = (4,3)
  \draw[-{Latex[length=3mm]}, thick] (0,0) -- (4,3);
  \node[above left=3pt] at (2.4,1.85) {$\vec v$};
  \node[right=8pt] at (4,3) {$(4,3)$};

  % Vector (-3,4)
  \draw[-{Latex[length=3mm]}, thick] (0,0) -- (-3,4);
  \node[left=8pt] at (-3,4) {$\begin{bmatrix}-3\\4\end{bmatrix} = -\vec{w}$};

  % Vector w = (3,-4)
  \draw[-{Latex[length=3mm]}, thick] (0,0) -- (3,-4);
  \node[right=8pt] at (3,-4) {$\begin{bmatrix}3\\-4\end{bmatrix}=\vec w$};

  % Origin
  \filldraw[black] (0,0) circle (2pt);
\end{tikzpicture}
\caption{The coefficient vector $\begin{bmatrix}3\\-4\end{bmatrix}$ is perpendicular to the line $3x-4y=0$.}
\end{figure}

Notice that the coefficients of the equation defining $\ell$ are

$$
(3,-4).
$$

Let us name the corresponding vector

$$
\begin{bmatrix}
3\\
-4
\end{bmatrix}
=
\vec w.
$$

Notice that $\vec w $
lies on $\ell^\perp$.
What is going on?

We can write

$$
3x-4y=0
$$

as

$$
\begin{bmatrix}
3\\
-4
\end{bmatrix}
\cdot
\begin{bmatrix}
x\\
y
\end{bmatrix}
=0.
$$

So the points on the line $\ell$ are exactly those that are perpendicular to
the vector

$$
\begin{bmatrix}
3\\
-4
\end{bmatrix}.
$$

\bigidea{
This is a rather remarkable perspective.

The coefficients of the equation defining $\ell$ are exactly given by vectors
on $\ell^\perp$.}

So we can think of the line $\ell$ in two different ways:

$$
\ell
=
\operatorname{span}(\vec v),
\qquad
\vec v=
\begin{bmatrix}
4\\
3
\end{bmatrix},
$$

or

$$
\ell
=
\left\{
\text{vectors perpendicular to } \vec w
\right\},
\qquad
\vec w=
\begin{bmatrix}
3\\
-4
\end{bmatrix}.
$$

Likewise, we can think of $\ell^\perp$ as

$$
\ell^\perp
=
\operatorname{span}(\vec w),
\qquad
\vec w=
\begin{bmatrix}
3\\
-4
\end{bmatrix},
$$

or

$$
\ell^\perp
=
\left\{
\text{vectors perpendicular to } \vec v
\right\},
\qquad
\vec v=
\begin{bmatrix}
4\\
3
\end{bmatrix}.
$$

So an equation for $\ell^\perp$ is

$$
4x+3y=0.
$$

\begin{remark}
    We can pick other vectors that span $\ell^\perp$ eg. 
    $$
    \ell^\perp = \operatorname{span} (3\vec{w})$$,
    where 
    $$ 3\vec{w} =\begin{bmatrix}
9\\
-12
\end{bmatrix} $$
This means another way of describing
$\ell$ is the equation
$$ 9x - 12y=0$$.
This just corresponds to multiplying the original equation we had by the scalar $3$.
\end{remark}

\paragraph{Exercise.}
Consider the line

$$
\ell'=\operatorname{span}
\begin{bmatrix}
2\\
1
\end{bmatrix}.
$$

\begin{enumerate}
  \item Draw $\ell'$ and $(\ell')^\perp$.
  \item Find a nonzero vector on $(\ell')^\perp$.
  \item Use your vector from the previous subpart to write an equation for $\ell'$.
  \item Find another nonzero vector on $(\ell')^\perp$.
  \item Use your vector from the previous subpart to write another equation for $\ell'$.
  \item How do your equations for $\ell'$ differ?
  \item Write down an equation for $\ell'$.
\end{enumerate}

\subsection{Affine Lines}

Now we can discuss lines that don't necessarily pass through the origin. These
are often called affine lines to indicate that they need not pass through the
origin.

How to describe them?

\begin{figure}[h]
\centering
\begin{tikzpicture}[scale=0.72]
  % Axes
  \draw[->, gray] (-5.5,0) -- (6.5,0) node[right] {$x$};
  \draw[->, gray] (0,-4.5) -- (0,4.8) node[above] {$y$};

  % Homogeneous line 3x - 4y = 0
  \draw[black, thick, domain=-5.2:5.2] plot (\x,{0.75*\x});
 

  \draw[-{Latex[length=3mm]}, ultra thick] (0,0) -- (4,3);
  \node[above left=3pt] at (2.4,1.85) {$\vec v$};

  % Affine line 3x - 4y = 7
  \draw[magenta!80!black, thick, domain=-1.2:6.2] plot (\x,{(3*\x - 7)/4});

  % Points on homogeneous line
  \filldraw[black] (-4,-3) circle (2pt);
  \node[below left] at (-4,-3) {$(-4,-3)$};

  \filldraw[black] (4,3) circle (2pt);
  \node[right=8pt] at (4,3) {$(4,3)$};

  % Points on affine line
  \filldraw[magenta!80!black] (1,-1) circle (2pt);
  \node[below right=3pt, magenta!80!black] at (1,-1) {$(1,-1)$};

  \filldraw[magenta!80!black] (5,2) circle (2pt);
  \node[right=8pt, magenta!80!black] at (5,2) {$(5,2)$};

  % Arrow along affine line
  \draw[-{Latex[length=3mm]}, magenta!80!black, ultra thick] (1,-1) -- (5,2);

  % Labels
  \node[violet, above left] at (5.05,4.05) {$\ell: 3x-4y=0$};
  

  \node[magenta!80!black, above right] at (6.2,2.75) {$\ell': 3x-4y=7$};
  
\end{tikzpicture}
\caption{The affine line $3x-4y=7$ is parallel to the homogeneous line $3x-4y=0$.}
\end{figure}

Let's consider the line

$$
3x-4y=7.
$$

Notice the point

$$
(1,-1)
$$

lies on it.

We can get any point on this line by first going to

$$
\begin{bmatrix}
1\\
-1
\end{bmatrix}
$$

then adding any vector on $\ell$ to it.

For this line, we may take

$$
\vec v=
\begin{bmatrix}
4\\
3
\end{bmatrix},
$$

since this vector points in the same direction as the line $3x-4y=0$.

So the line $3x-4y=7$ is

$$
\begin{bmatrix}
1\\
-1
\end{bmatrix}
+
\operatorname{span}(\vec v)
=
\left\{
\begin{bmatrix}
1\\
-1
\end{bmatrix}
+
t\vec v
\;:\;
t\in\mathbb{R}
\right\}.
$$

Instead of

$$
\begin{bmatrix}
1\\
-1
\end{bmatrix},
$$

we could have taken any point on the line, e.g.

$$
\begin{bmatrix}
5\\
2
\end{bmatrix},
$$

so

$$
3x-4y=7
$$

can also be written as

$$
\left\{
\begin{bmatrix}
5\\
2
\end{bmatrix}
+
t\vec v'
\;:\;
t\in\mathbb{R}
\right\}.
$$

\paragraph{Summary.}
Affine lines in $\mathbb{R}^2$ may be thought of as translations of the span
of a vector by some fixed vector.

\paragraph{Exercise.}
Can you describe affine lines using the dual perspective and orthogonality?

\paragraph{Answer.}
They can be described as a set of vectors $\vec{v'}$ such that

$$
\vec{v'}\cdot \vec w
$$

is a fixed constant, where $\vec w$ is a nonzero vector on $\ell^\perp$.

For example, the affine line

$$
3x-4y=7
$$

can be written as

$$
\begin{bmatrix}
3\\
-4
\end{bmatrix}
\cdot
\begin{bmatrix}
x\\
y
\end{bmatrix}
=7.
$$

So it is the set of all points whose dot product with the fixed normal vector

$$
\begin{bmatrix}
3\\
-4
\end{bmatrix}
$$

is the fixed constant $7$.


% Requires:
% \usepackage{amsmath,amssymb}
% \usepackage{tikz}
% \usetikzlibrary{arrows.meta,positioning}


\section{Geometry in \texorpdfstring{$\mathbb{R}^2$}{R2}}

There are two special vectors in $\mathbb{R}^2$,

\[
\vec e_1=
\begin{bmatrix}
1\\
0
\end{bmatrix},
\qquad
\vec e_2=
\begin{bmatrix}
0\\
1
\end{bmatrix}.
\]

What is nice about them is that every vector can be easily expressed as a
combination of them.

\begin{figure}[h]
\centering
\begin{tikzpicture}[scale=1.0]
  % Axes
  \draw[->, gray] (-2.0,0) -- (3.5,0) node[right] {$x$};
  \draw[->, gray] (0,-1.5) -- (0,3.5) node[above] {$y$};

  % e1
  \draw[-{Latex[length=3mm]}, thick] (0,0) -- (2.0,0);
  \node[below right] at (2.0,0) {$\vec e_1=\begin{bmatrix}1\\0\end{bmatrix}$};

  % e2
  \draw[-{Latex[length=3mm]}, thick] (0,0) -- (0,2.0);
  \node[above left] at (0,2.0) {$\vec e_2=\begin{bmatrix}0\\1\end{bmatrix}$};

  \filldraw[black] (0,0) circle (1.5pt);
\end{tikzpicture}
\caption{The standard basis vectors in $\mathbb{R}^2$.}
\end{figure}

For example,

\[
\begin{bmatrix}
5\\
7
\end{bmatrix}
=
\begin{bmatrix}
5\\
0
\end{bmatrix}
+
\begin{bmatrix}
0\\
7
\end{bmatrix}
=
5
\begin{bmatrix}
1\\
0
\end{bmatrix}
+
7
\begin{bmatrix}
0\\
1
\end{bmatrix}
=
5\vec e_1+7\vec e_2.
\]

To do geometry in $\mathbb{R}^2$, it is useful to consider other such pairs of
vectors. The key properties of $(\vec e_1,\vec e_2)$ are:

\begin{enumerate}
  \item length one,
  \item orthogonal to each other.
\end{enumerate}

We will now discuss the notion of an orthonormal basis.

For this, we first need to define unit vectors.

\paragraph{Definition.}
A unit vector is a vector of length one.

\paragraph{Definition.}
An orthonormal basis of $\mathbb{R}^2$ is a pair of vectors
$(\vec u_1,\vec u_2)$ such that $\vec u_1$ and $\vec u_2$ are both unit vectors
that are orthogonal to each other.

\paragraph{Example.}
The vector

\[
\begin{bmatrix}
\frac45\\[2pt]
\frac35
\end{bmatrix}
\]

has length one since

\[
\sqrt{\left(\frac45\right)^2+\left(\frac35\right)^2}
=
\sqrt{\frac{16}{25}+\frac{9}{25}}
=
1.
\]

Similarly, the vector

\[
\begin{bmatrix}
\frac{1}{\sqrt{2}}\\[2pt]
-\frac{1}{\sqrt{2}}
\end{bmatrix}
\]

has length one since

\[
\sqrt{
\left(\frac{1}{\sqrt{2}}\right)^2+
\left(-\frac{1}{\sqrt{2}}\right)^2
}
=
\sqrt{\frac12+\frac12}
=
1.
\]

\begin{figure}[h]
\centering
\begin{tikzpicture}[scale=1.3]
  % Axes
  \draw[->, gray] (-1.6,0) -- (2.2,0) node[right] {$x$};
  \draw[->, gray] (0,-1.6) -- (0,2.0) node[above] {$y$};

  % Vector (4/5,3/5)
  \draw[-{Latex[length=3mm]}, thick] (0,0) -- (1.6,1.2);
  \node[right] at (1.6,1.2)
    {$\begin{bmatrix}\frac45\\[2pt]\frac35\end{bmatrix}$};

  % Vector (1/sqrt2,-1/sqrt2)
  \draw[-{Latex[length=3mm]}, thick] (0,0) -- (1.35,-1.35);
  \node[right] at (1.35,-1.35)
    {$\begin{bmatrix}\frac{1}{\sqrt{2}}\\[2pt]-\frac{1}{\sqrt{2}}\end{bmatrix}$};

  \filldraw[black] (0,0) circle (1.5pt);
\end{tikzpicture}
\caption{Examples of unit vectors in $\mathbb{R}^2$.}
\end{figure}

\paragraph{Example.}
Let

\[
\vec u_1=
\begin{bmatrix}
\frac45\\[2pt]
\frac35
\end{bmatrix},
\qquad
\vec u_2=
\begin{bmatrix}
-\frac35\\[2pt]
\frac45
\end{bmatrix}.
\]

Then $(\vec u_1,\vec u_2)$ is an orthonormal basis of $\mathbb{R}^2$.

\begin{figure}[h]
\centering
\begin{tikzpicture}[scale=1.3]
  % Axes
  \draw[->, gray] (-2.0,0) -- (2.3,0) node[right] {$x$};
  \draw[->, gray] (0,-1.5) -- (0,2.3) node[above] {$y$};

  % u1
  \draw[-{Latex[length=3mm]}, thick] (0,0) -- (1.6,1.2);
  \node[right] at (1.6,1.2)
    {$\vec u_1=\begin{bmatrix}\frac45\\[2pt]\frac35\end{bmatrix}$};

  % u2
  \draw[-{Latex[length=3mm]}, thick] (0,0) -- (-1.2,1.6);
  \node[left] at (-1.2,1.6)
    {$\vec u_2=\begin{bmatrix}-\frac35\\[2pt]\frac45\end{bmatrix}$};

  \filldraw[black] (0,0) circle (1.5pt);
\end{tikzpicture}
\caption{The pair $(\vec u_1,\vec u_2)$ is an orthonormal basis of $\mathbb{R}^2$.}
\end{figure}

\paragraph{Example.}
Let

\[
\vec w_1=
\begin{bmatrix}
-\frac{1}{\sqrt{2}}\\[2pt]
\frac{1}{\sqrt{2}}
\end{bmatrix},
\qquad
\vec w_2=
\begin{bmatrix}
-\frac{1}{\sqrt{2}}\\[2pt]
-\frac{1}{\sqrt{2}}
\end{bmatrix}.
\]

Then $(\vec w_1,\vec w_2)$ is also an orthonormal basis.

\begin{figure}[h]
\centering
\begin{tikzpicture}[scale=1.4]
  % Axes
  \draw[->, gray] (-2.2,0) -- (1.8,0) node[right] {$x$};
  \draw[->, gray] (0,-1.9) -- (0,2.0) node[above] {$y$};

  % w1
  \draw[-{Latex[length=3mm]}, thick] (0,0) -- (-1.25,1.25);
  \node[left] at (-1.25,1.25)
    {$\vec w_1=\begin{bmatrix}-\frac{1}{\sqrt{2}}\\[2pt]\frac{1}{\sqrt{2}}\end{bmatrix}$};

  % w2
  \draw[-{Latex[length=3mm]}, thick] (0,0) -- (-1.25,-1.25);
  \node[left] at (-1.25,-1.25)
    {$\vec w_2=\begin{bmatrix}-\frac{1}{\sqrt{2}}\\[2pt]-\frac{1}{\sqrt{2}}\end{bmatrix}$};

  \filldraw[black] (0,0) circle (1.5pt);
\end{tikzpicture}
\caption{An orthonormal basis is not unique.}
\end{figure}

\paragraph{Remark.}
Note that an orthonormal basis is not unique.

\subsection{Why Are Orthonormal Bases Important?}

The reason orthonormal bases are useful is that it is very easy to express an
arbitrary vector in terms of an orthonormal basis.

Let $\vec u_1,\vec u_2$ be an orthonormal basis of $\mathbb{R}^2$.

Let $\vec v$ be any vector in $\mathbb{R}^2$.

Then

\[
\vec v=a\vec u_1+b\vec u_2
\]

for some scalars $a,b$. This follows by completing the rectangle in the figure below. 

But how do we find $a,b$?

Take the dot product with $\vec u_1$:

\[
\vec v\cdot \vec u_1
=
(a\vec u_1+b\vec u_2)\cdot \vec u_1
=
a(\vec u_1\cdot \vec u_1)
+
b(\vec u_2\cdot \vec u_1).
\]

Since $\vec u_1$ is a unit vector,

\[
\vec u_1\cdot \vec u_1=1.
\]

Since $\vec u_1$ and $\vec u_2$ are orthogonal,

\[
\vec u_2\cdot \vec u_1=0.
\]

Therefore

\[
\vec v\cdot \vec u_1=a.
\]

Similarly,

\[
b=\vec v\cdot \vec u_2.
\]

So

\[
\vec v
=
(\vec v\cdot \vec u_1)\vec u_1
+
(\vec v\cdot \vec u_2)\vec u_2.
\]



\begin{figure}[h]
\centering
\begin{tikzpicture}[scale=1.0]
  % Orthonormal unit vectors:
  % u_1 = (3/sqrt(10), 1/sqrt(10))
  % u_2 = (-1/sqrt(10), 3/sqrt(10))
  %
  % Components:
  % a u_1 = 3u_1
  % b u_2 = (5/2)u_2
  % v = a u_1 + b u_2

  \coordinate (O) at (0,0);

  \coordinate (Uone) at (0.94868,0.31623);
  \coordinate (Utwo) at (-0.31623,0.94868);

  \coordinate (A) at (2.84605,0.94868);
  \coordinate (B) at (-0.79057,2.37171);

  % V = A + B
  \coordinate (V) at (2.05548,3.32039);

  % Axes
  \draw[->, gray] (-2.0,0) -- (4.2,0) node[right] {$x$};
  \draw[->, gray] (0,-1.5) -- (0,4.2) node[above] {$y$};

  % Dashed parallelogram guides
  \draw[dashed, gray] (A) -- (V);
  \draw[dashed, gray] (B) -- (V);

  % Component a u_1
  \draw[-{Latex[length=3mm]}, thick] (O) -- (A);
  \node[below right=3pt] at (A) {$a\vec u_1$};

  % Component b u_2
  \draw[-{Latex[length=3mm]}, thick] (O) -- (B);
  \node[left=5pt] at (B) {$b\vec u_2$};

  % Unit vector u_1
  \draw[-{Latex[length=2.2mm]}, thick] (O) -- (Uone);
  \node[above=3pt] at (Uone) {$\vec u_1$};

  % Unit vector u_2
  \draw[-{Latex[length=2.2mm]}, thick] (O) -- (Utwo);
  \node[right=3pt] at (Utwo) {$\vec u_2$};

  % Vector v = a u_1 + b u_2
  \draw[-{Latex[length=3mm]}, very thick] (O) -- (V);
  \node[above right=3pt] at (V) {$\vec v$};

  % Origin
  \filldraw[black] (O) circle (1.5pt);
\end{tikzpicture}
\caption{The decomposition
$\vec v=a\vec u_1+b\vec u_2$ with respect to the orthonormal directions
$\vec u_1$ and $\vec u_2$.}
\end{figure}

\paragraph{Example.}
Let

\[
\vec v=
\begin{bmatrix}
7\\
5
\end{bmatrix}.
\]

For the standard basis

\[
(\vec e_1,\vec e_2)
=
\left(
\begin{bmatrix}
1\\
0
\end{bmatrix},
\begin{bmatrix}
0\\
1
\end{bmatrix}
\right),
\]

we have

\[
\vec v\cdot \vec e_1
=
\begin{bmatrix}
7\\
5
\end{bmatrix}
\cdot
\begin{bmatrix}
1\\
0
\end{bmatrix}
=
7,
\]

and

\[
\vec v\cdot \vec e_2
=
\begin{bmatrix}
7\\
5
\end{bmatrix}
\cdot
\begin{bmatrix}
0\\
1
\end{bmatrix}
=
5.
\]

Therefore

\[
(\vec v\cdot \vec e_1)\vec e_1
+
(\vec v\cdot \vec e_2)\vec e_2
=
7\vec e_1+5\vec e_2
=
\begin{bmatrix}
7\\
0
\end{bmatrix}
+
\begin{bmatrix}
0\\
5
\end{bmatrix}
=
\begin{bmatrix}
7\\
5
\end{bmatrix}.
\]

Now take the orthonormal basis

\[
(\vec u_1,\vec u_2)
=
\left(
\begin{bmatrix}
\frac45\\[2pt]
\frac35
\end{bmatrix},
\begin{bmatrix}
-\frac35\\[2pt]
\frac45
\end{bmatrix}
\right).
\]

Then

\[
\vec v\cdot \vec u_1
=
\begin{bmatrix}
7\\
5
\end{bmatrix}
\cdot
\begin{bmatrix}
\frac45\\[2pt]
\frac35
\end{bmatrix}
=
\frac{28}{5}+\frac{15}{5}
=
\frac{43}{5}.
\]

Also,

\[
\vec v\cdot \vec u_2
=
\begin{bmatrix}
7\\
5
\end{bmatrix}
\cdot
\begin{bmatrix}
-\frac35\\[2pt]
\frac45
\end{bmatrix}
=
-\frac{21}{5}+\frac{20}{5}
=
-\frac15.
\]

So

\[
\vec v
=
\frac{43}{5}\vec u_1
-
\frac15 \vec u_2.
\]


Continuing the example, suppose

\[
(\vec w_1,\vec w_2)
=
\left(
\begin{bmatrix}
-\frac{1}{\sqrt{2}}\\[2pt]
\frac{1}{\sqrt{2}}
\end{bmatrix},
\begin{bmatrix}
-\frac{1}{\sqrt{2}}\\[2pt]
-\frac{1}{\sqrt{2}}
\end{bmatrix}
\right).
\]

Let

\[
\vec v=
\begin{bmatrix}
7\\
5
\end{bmatrix}.
\]

Then

\[
\vec v\cdot \vec w_1
=
\begin{bmatrix}
7\\
5
\end{bmatrix}
\cdot
\begin{bmatrix}
-\frac{1}{\sqrt{2}}\\[2pt]
\frac{1}{\sqrt{2}}
\end{bmatrix}
=
-\frac{7}{\sqrt{2}}+\frac{5}{\sqrt{2}}
=
-\frac{2}{\sqrt{2}}
=
-\sqrt{2}.
\]

Also,

\[
\vec v\cdot \vec w_2
=
\begin{bmatrix}
7\\
5
\end{bmatrix}
\cdot
\begin{bmatrix}
-\frac{1}{\sqrt{2}}\\[2pt]
-\frac{1}{\sqrt{2}}
\end{bmatrix}
=
-\frac{7}{\sqrt{2}}-\frac{5}{\sqrt{2}}
=
-\frac{12}{\sqrt{2}}
=
-6\sqrt{2}.
\]

So

\[
\vec v
=
-\sqrt{2}\,\vec w_1
-
6\sqrt{2}\,\vec w_2.
\]

\subsection{Orthogonal Projection Onto a Line}


\begin{figure}[h]
\centering
\begin{tikzpicture}[scale=0.9]
  % Coordinates
  \coordinate (O) at (0,0);
  \coordinate (Vone) at (2.55,1.40);
  \coordinate (Vtwo) at (-1.25,2.27);
  \coordinate (V) at (1.30,3.67); % Vone + Vtwo

  % Axes
  \draw[->, gray] (-3.8,0) -- (4.2,0) node[right] {$x$};
  \draw[->, gray] (0,-3.2) -- (0,4.3) node[above] {$y$};

  % Line ell
  \draw[black, thick, domain=-3.7:4.0] plot (\x,{0.55*\x});
  \node[right] at (3.75,1.95) {$\ell$};

  % Perpendicular line ell perp
  \draw[black, thick, domain=-1.7:1.7] plot (\x,{-1.818*\x});
  \node[left] at (-1.45,3.4) {$\ell^\perp$};

  % Unit-ish arrow along ell
  \draw[-{Latex[length=2.5mm]}, thick] (O) -- (1.55,0.85);
  \node[below right] at (1.35,0.65) {$\vec u_1$};

  % Projection/vector along ell
  \draw[-{Latex[length=2.5mm]}, thick] (O) -- (Vone);
  \node[above] at (2.25,1.25) {$\vec v_1$};

  % Unit-ish arrow along ell perp
  \draw[-{Latex[length=2.5mm]}, thick] (O) -- (-0.65,1.18);
  \node[left] at (-0.65,1.18) {$\vec u_2$};

  % Projection/vector along ell perp
  \draw[-{Latex[length=2.5mm]}, thick] (O) -- (Vtwo);
  \node[left] at (-1.25,2.27) {$\vec v_2$};

  % Vector v = v1 + v2
  \draw[-{Latex[length=3mm]}, thick] (O) -- (V);
  \node[above right] at (V) {$\vec v$};

  % Dashed parallelogram guide lines
  \draw[dashed] (Vone) -- (V);
  \draw[dashed] (Vtwo) -- (V);

  % Origin
  \filldraw[black] (O) circle (1.5pt);
\end{tikzpicture}
\caption{A vector can be resolved into a part along a line and a part perpendicular to the line.}
\end{figure}


Now suppose $\vec v$ is some vector in $\mathbb{R}^2$, and $\ell$ is a line in
$\mathbb{R}^2$ going through the origin.

We can write $\vec v$ as the sum of two vectors, one along $\ell$, and one
perpendicular to $\ell$:

\[
\vec v=\vec v_1+\vec v_2.
\]

The vector $\vec v_1$ is called the orthogonal projection of $\vec v$ to
$\ell$.

How do we find it?

First, we can pick a unit vector $\vec u_1$ along $\ell$, and a unit vector
$\vec u_2$ perpendicular to $\ell$.

Then $(\vec u_1,\vec u_2)$ is an orthonormal basis of $\mathbb{R}^2$.

We saw that

\[
\vec v
=
(\vec v\cdot \vec u_1)\vec u_1
+
(\vec v\cdot \vec u_2)\vec u_2.
\]

The first term is along $\ell$, and the second term is perpendicular to $\ell$.
Therefore

\[
\operatorname{proj}_{\ell}(\vec v)
=
\vec v_1
=
(\vec v\cdot \vec u_1)\vec u_1.
\]

\paragraph{Example.}
Find the orthogonal projection of the vector

\[
\begin{bmatrix}
8\\
9
\end{bmatrix}
\]

to the line

\[
2x-3y=0.
\]

\begin{figure}[h]
\centering
\begin{tikzpicture}[scale=0.55]
  % Axes
  \draw[->, gray] (-3.5,0) -- (10.5,0) node[right] {$x$};
  \draw[->, gray] (0,-3.5) -- (0,10.5) node[above] {$y$};

  % Line ell: 2x - 3y = 0 => y = 2x/3
  \draw[black, thick, domain=-3:10] plot (\x,{(2/3)*\x});
  \node[right] at (8.7,5.63) {$\ell$};
  \node[right] at (7.0,4.2) {$2x-3y=0$};

  % Vector v = (8,9)
  \draw[-{Latex[length=3mm]}, thick] (0,0) -- (8,9);
  \node[above right] at (8,9)
    {$\vec v=\begin{bmatrix}8\\9\end{bmatrix}$};

  % Unit direction u1, scaled visually
  \draw[-{Latex[length=3mm]}, thick] (0,0) -- (3,2);
  \node[below right] at (3,2) {$\vec u_1$};

  % Projection of v onto ell = (126/13,84/13)
  \draw[-{Latex[length=3mm]}, thick] (0,0) -- ({126/13},{84/13});
  \node[below right] at ({126/13},{84/13})
    {$\operatorname{proj}_{\ell}(\vec v)$};

  % Perpendicular component
  \draw[dashed] ({126/13},{84/13}) -- (8,9);

  \filldraw[black] (0,0) circle (2pt);
\end{tikzpicture}
\caption{Projection of $\vec v=\begin{bmatrix}8\\9\end{bmatrix}$ onto the line $2x-3y=0$.}
\end{figure}

A vector along the line $2x-3y=0$ is

\[
\begin{bmatrix}
3\\
2
\end{bmatrix}
=
\vec w.
\]

Its length is

\[
\lVert \vec w\rVert
=
\sqrt{9+4}
=
\sqrt{13}.
\]

A unit vector along $\ell$ is

\[
\vec u_1
=
\begin{bmatrix}
\frac{3}{\sqrt{13}}\\[2pt]
\frac{2}{\sqrt{13}}
\end{bmatrix}.
\]

So the projection is

\[
\operatorname{proj}_{\ell}(\vec v)
=
(\vec v\cdot \vec u_1)\vec u_1.
\]

Thus

\[
\operatorname{proj}_{\ell}(\vec v)
=
\left(
\begin{bmatrix}
8\\
9
\end{bmatrix}
\cdot
\begin{bmatrix}
\frac{3}{\sqrt{13}}\\[2pt]
\frac{2}{\sqrt{13}}
\end{bmatrix}
\right)
\begin{bmatrix}
\frac{3}{\sqrt{13}}\\[2pt]
\frac{2}{\sqrt{13}}
\end{bmatrix}.
\]

Since

\[
\begin{bmatrix}
8\\
9
\end{bmatrix}
\cdot
\begin{bmatrix}
\frac{3}{\sqrt{13}}\\[2pt]
\frac{2}{\sqrt{13}}
\end{bmatrix}
=
\frac{24}{\sqrt{13}}+\frac{18}{\sqrt{13}}
=
\frac{42}{\sqrt{13}},
\]

we get

\[
\operatorname{proj}_{\ell}(\vec v)
=
\frac{42}{\sqrt{13}}
\begin{bmatrix}
\frac{3}{\sqrt{13}}\\[2pt]
\frac{2}{\sqrt{13}}
\end{bmatrix}
=
\begin{bmatrix}
\frac{126}{13}\\[2pt]
\frac{84}{13}
\end{bmatrix}.
\]

Likewise, we can find the projection to $\ell^\perp$.

Since

\[
\vec u_1
=
\begin{bmatrix}
\frac{3}{\sqrt{13}}\\[2pt]
\frac{2}{\sqrt{13}}
\end{bmatrix},
\]

we can take

\[
\vec u_2
=
\begin{bmatrix}
-\frac{2}{\sqrt{13}}\\[2pt]
\frac{3}{\sqrt{13}}
\end{bmatrix}.
\]

Then

\[
\operatorname{proj}_{\ell^\perp}(\vec v)
=
(\vec v\cdot \vec u_2)\vec u_2.
\]

So

\[
\operatorname{proj}_{\ell^\perp}(\vec v)
=
\left(
\begin{bmatrix}
8\\
9
\end{bmatrix}
\cdot
\begin{bmatrix}
-\frac{2}{\sqrt{13}}\\[2pt]
\frac{3}{\sqrt{13}}
\end{bmatrix}
\right)
\begin{bmatrix}
-\frac{2}{\sqrt{13}}\\[2pt]
\frac{3}{\sqrt{13}}
\end{bmatrix}.
\]

Since

\[
\begin{bmatrix}
8\\
9
\end{bmatrix}
\cdot
\begin{bmatrix}
-\frac{2}{\sqrt{13}}\\[2pt]
\frac{3}{\sqrt{13}}
\end{bmatrix}
=
-\frac{16}{\sqrt{13}}+\frac{27}{\sqrt{13}}
=
\frac{11}{\sqrt{13}},
\]

we get

\[
\operatorname{proj}_{\ell^\perp}(\vec v)
=
\frac{11}{\sqrt{13}}
\begin{bmatrix}
-\frac{2}{\sqrt{13}}\\[2pt]
\frac{3}{\sqrt{13}}
\end{bmatrix}
=
\begin{bmatrix}
-\frac{22}{13}\\[2pt]
\frac{33}{13}
\end{bmatrix}.
\]

Then

\[
\vec v
=
\vec v_1+\vec v_2,
\]

where

\[
\vec v_1=
\begin{bmatrix}
\frac{126}{13}\\[2pt]
\frac{84}{13}
\end{bmatrix},
\qquad
\vec v_2=
\begin{bmatrix}
-\frac{22}{13}\\[2pt]
\frac{33}{13}
\end{bmatrix}.
\]

Check:

\[
\begin{bmatrix}
\frac{126}{13}\\[2pt]
\frac{84}{13}
\end{bmatrix}
+
\begin{bmatrix}
-\frac{22}{13}\\[2pt]
\frac{33}{13}
\end{bmatrix}
=
\begin{bmatrix}
\frac{104}{13}\\[2pt]
\frac{117}{13}
\end{bmatrix}
=
\begin{bmatrix}
8\\
9
\end{bmatrix}.
\]

Basically, we have resolved one vector $\vec v$ into two orthogonal directions,
with $\vec v_1$ along $\ell$ and $\vec v_2$ along $\ell^\perp$.
